% Eigene Befehle und typographische Auszeichnungen für diese

% einfaches Wechseln der Schrift, z.B.: \changefont{cmss}{sbc}{n}
\newcommand{\changefont}[3]{\fontfamily{#1} \fontseries{#2} \fontshape{#3} \selectfont}

% Abkürzungen mit korrektem Leerraum 
\newcommand{\ua}{\mbox{u.\,a.\ }}
\newcommand{\zB}{\mbox{z.\,B.\ }}
\newcommand{\dahe}{\mbox{d.\,h.\ }}
\newcommand{\Vgl}{Vgl.\ }
\newcommand{\bzw}{bzw.\ }
\newcommand{\evtl}{evtl.\ }
\newcommand{\ie}{\mbox{i.\,e.\ }}

\newcommand{\illustration}[1]{illustration~\ref{fig:#1}}

\newcommand{\bs}{$\backslash$}

\newcommand{\sectiontitle}{}
\let\TeXsection\section
\renewcommand{\section}[2][\empty]{\ifthenelse{\equal{#1}{\empty}}{\TeXsection{#2}\renewcommand{\sectiontitle}{#2}}{\TeXsection{#2}\label{#1}\renewcommand{\sectiontitle}{#1}}}

\newcommand{\subsectiontitle}{}
\let\TeXsubsection\subsection
\renewcommand{\subsection}[2][\empty]{\ifthenelse{\equal{#1}{\empty}}{\TeXsubsection{#2}\renewcommand{\subsectiontitle}{#2}}{\TeXsubsection{#2}\label{#1}\renewcommand{\subsectiontitle}{#1}}}

\let\TeXsubsubsection\subsubsection
\renewcommand{\subsubsection}[2][\empty]{\ifthenelse{\equal{#1}{\empty}}{\TeXsubsubsection{#2}}{\TeXsubsubsection{#2}\label{#1}}}

\newcommand{\supref}[1]{\textsuperscript{\ref{#1}}}

% erzeugt ein Listenelement mit fetter Überschrift 
\newcommand{\itemd}[2]{\item{\textbf{#1}}\\{#2}}

% einige Befehle zum Zitieren --------------------------------------------------
\newcommand{\Zitat}[2][\empty]{\ifthenelse{\equal{#1}{\empty}}{\citep{#2}}{\citep[#1]{#2}}}

% zum Ausgeben von Autoren
\newcommand{\AutorName}[1]{\textsc{#1}}
\newcommand{\Autor}[1]{\AutorName{\citeauthor{#1}}}

% verschiedene Befehle um Wörter semantisch auszuzeichnen ----------------------
\newcommand{\newDef}[1]{\textbf{#1}}
\newcommand{\techTerm}[1]{\textit{#1}}

\newcommand{\DisplInput}[1]{\texttt{#1}}
\newcommand{\DisplCode}[1]{\texttt{#1}}
\newcommand{\DisplFile}[1]{\texttt{#1}}

\newcommand{\Datentyp}[1]{\textsf{#1}}
\newcommand{\XMLElement}[1]{\textsf{#1}}
\newcommand{\Webservice}[1]{\textsf{#1}}

% Deklaration von Hinweistexten
\newcommand{\hint}[1]{
\begin{tabular}{w{1.5cm}p{12cm}}\\
\scalebox{0.125}{\includegraphics{images/base/Hint_Notification}} & #1
\end{tabular}
}

% Deklaration von Warnungen
\newcommand{\warning}[1]{
\begin{tabular}{w{1.5cm}p{12cm}}\\
\scalebox{0.125}{\includegraphics{images/base/Warning_Notification}} & #1
\end{tabular}
}

% einbinden eines emp Bildes.
\newcommand{\useempcenter}[2]{\begin{figure}[h!]\centering\scalebox{0.8}{\empuse{#1}}\caption{#2}\end{figure}}

\newcommand{\useempfloat}[3]{\begin{wrapfigure}{#1}{0.4\textwidth}\begin{center}\empuse{#2}\end{center}\caption{#3}\end{wrapfigure}}

\newcommand{\useemp}[3][\empty]{\ifthenelse{\equal{#1}{\empty}}{\useempcenter{#2}{#3}}{\useempfloat{#1}{#2}{#3}}}

% Wrapper für useemp für uml Klassendiagramme
\newcommand{\umlclass}[3][\empty]{\ifthenelse{\equal{#1}{\empty}}{\useemp{#2}{#3}}{\useemp[#1]{#2}{#3}}}

% Wrapper um listings einzubinden
\newcommand{\phplisting}[3]{\lstset{language=PHP, basicstyle=\footnotesize, showstringspaces=false, tabsize=2}\lstinputlisting[label=#1,caption=#2]{#3}}
\newcommand{\xmllisting}[3]{\lstset{language=XML, basicstyle=\footnotesize, showstringspaces=false, tabsize=2}\lstinputlisting[label=#1,caption=#2]{#3}}
